% ===========================
% Navier–Stokes Global Regularity via Categorical Collapse
% ===========================
\documentclass[11pt]{article}

% === Language and Encoding ===
\usepackage[utf8]{inputenc}
\usepackage[T1]{fontenc}
\usepackage[english]{babel}

% === Math and Symbols ===
\usepackage{amsmath,amssymb,amsthm,amsfonts}
\usepackage{mathtools}
\usepackage{mathrsfs}
\usepackage{stmaryrd}  % for \llbracket etc.
\usepackage{bm}        % bold math
\usepackage{tikz}
\usetikzlibrary{matrix, arrows.meta, calc, decorations.pathmorphing}

% === Code, Listings, Diagrams ===
\usepackage{listings}
\lstset{
  basicstyle=\ttfamily\footnotesize,
  keywordstyle=\color{blue},
  commentstyle=\color{gray},
  breaklines=true,
  frame=single,
  captionpos=b
}

% === Graphics and Layout ===
\usepackage{graphicx}
\usepackage{xcolor}
\usepackage{geometry}
\geometry{margin=1in}

% === Hyperlinks ===
\usepackage[colorlinks=true, linkcolor=blue, citecolor=blue, urlcolor=blue]{hyperref}

% === Title Metadata ===
\title{Navier–Stokes Global Regularity via Categorical Collapse\\
\Large \textsc{Version 6.0}\\
\small Based on the AK High-Dimensional Projection Structural Theory v10.0}
\author{Atsushi Kobayashi \\ \small with ChatGPT Research Partner}
\date{June 2025}

% === Theorem Environments ===
\newtheorem{theorem}{Theorem}[section]
\newtheorem{definition}[theorem]{Definition}
\newtheorem{lemma}[theorem]{Lemma}
\newtheorem{corollary}[theorem]{Corollary}
\newtheorem{proposition}[theorem]{Proposition}
\newtheorem{remark}[theorem]{Remark}
\newtheorem{example}[theorem]{Example}

% === Math Operators ===
\DeclareMathOperator{\Ext}{Ext}
\DeclareMathOperator{\Hom}{Hom}
\DeclareMathOperator{\Spec}{Spec}
\DeclareMathOperator{\colim}{colim}
\DeclareMathOperator{\PH}{PH}
\DeclareMathOperator{\Tor}{Tor}
\DeclareMathOperator{\rank}{rank}
\DeclareMathOperator{\im}{im}
\DeclareMathOperator{\id}{id}
\DeclareMathOperator{\Ker}{Ker}
\DeclareMathOperator{\Coker}{Coker}

% === Custom Shortcuts ===
\newcommand{\QQ}{\mathbb{Q}}
\newcommand{\RR}{\mathbb{R}}
\newcommand{\CC}{\mathbb{C}}
\newcommand{\ZZ}{\mathbb{Z}}
\newcommand{\TT}{\mathbb{T}}

\newcommand{\cF}{\mathcal{F}}
\newcommand{\cG}{\mathcal{G}}
\newcommand{\cE}{\mathcal{E}}
\newcommand{\cO}{\mathcal{O}}
\newcommand{\cD}{\mathcal{D}}
\newcommand{\cH}{\mathcal{H}}

\newcommand{\into}{\hookrightarrow}
\newcommand{\onto}{\twoheadrightarrow}
\newcommand{\eps}{\varepsilon}
\newcommand{\Sha}{\mathbb{S}}

% === Document Starts ===
\begin{document}

\maketitle
\tableofcontents
\newpage




\section{Chapter 1: Introduction and Problem Setup}

\subsection{1.1 The Navier--Stokes Equations}

The three-dimensional incompressible Navier--Stokes equations describe the motion of viscous fluid flows and are given by:
\[
\partial_t u + (u \cdot \nabla)u = -\nabla p + \nu \Delta u, \quad \nabla \cdot u = 0,
\]
where \( u(t,x) \in \mathbb{R}^3 \) is the velocity field, \( p(t,x) \) is the pressure, and \( \nu > 0 \) is the kinematic viscosity. These equations model a wide class of physical phenomena including air, water, and other classical fluids.

Despite their apparent simplicity, the nonlinear convective term \( (u \cdot \nabla)u \) creates profound mathematical challenges, particularly regarding the global-in-time behavior of solutions.

\subsection{1.2 The Millennium Problem: Global Regularity}

The Clay Mathematics Institute has identified the following fundamental open problem as one of the Millennium Prize Problems:

\begin{quote}
\textit{Given smooth initial data with finite energy, does there exist a unique, globally smooth solution \( u(t,x) \in C^\infty(\mathbb{R}^3 \times [0,\infty)) \) to the 3D incompressible Navier--Stokes equations, or can singularities develop in finite time?}
\end{quote}

This question remains unresolved despite substantial analytical progress. Known results confirm existence of global weak solutions (Leray–Hopf solutions), but their regularity and uniqueness remain elusive in three dimensions.

\subsection{1.3 Previous Approaches and Limitations}

Many analytic techniques have been developed to address this problem, including:

\begin{itemize}
  \item \textbf{Ladyzhenskaya–Prodi–Serrin criteria} based on integrability bounds of the velocity field,
  \item \textbf{Beale–Kato–Majda (BKM) criterion} linking blow-up to the growth of vorticity,
  \item \textbf{Energy methods} focusing on bounding the \( H^1 \) norm of solutions,
  \item \textbf{Geometric and spectral approaches} examining scale-invariance and energy cascades,
  \item \textbf{Numerical simulations} to trace singularity formation and vortex stretching.
\end{itemize}

However, no existing framework has succeeded in proving or disproving global regularity from a structurally complete, obstruction-free viewpoint that is both semantically rigorous and constructively realizable.

\subsection{1.4 Our Approach: Categorical Collapse Framework}

In this paper, we propose a novel approach grounded in a structural and categorical interpretation of the Navier--Stokes dynamics. Our central claim is the following:

\begin{quote}
\textit{Global regularity can be obtained by showing that topological and categorical obstructions collapse to triviality over time, and this collapse is provable in a type-theoretic framework.}
\end{quote}

Rather than bounding norms directly or controlling function spaces globally, we encode the evolution of the velocity field into a derived sheaf structure \( \mathcal{F}_t \), and track the persistent homological and Ext-class complexity of this sheaf.

Two collapse routes are then established:

\begin{enumerate}
  \item Vanishing of persistent homology \( \mathrm{PH}_1(\mathcal{F}_t) \Rightarrow u(t) \in C^\infty \),
  \item Triviality of Ext-class \( \mathrm{Ext}^1(\mathcal{F}_t, \mathbb{Q}) = 0 \Rightarrow u(t) \in C^\infty \).
\end{enumerate}

These routes correspond to the removal of topological loops and categorical gluing obstructions, respectively.

\subsection{1.5 Foundational Basis: AK Theory}

The theoretical foundation for this approach is the \textbf{AK High-Dimensional Projection Structural Theory (AK-HDPST)}, version 10.0. This framework introduces:

\begin{itemize}
  \item A categorical collapse functor mapping filtered sheaves to trivial objects,
  \item A collapse zone in which persistent homology and Ext-classes vanish,
  \item A type-theoretic encoding of smoothness conditions using dependent types,
  \item ZFC-consistent semantics linking geometric evolution to logical form.
\end{itemize}

Specifically, we make use of:

\begin{itemize}
  \item \textbf{Chapter 8} of AK-HDPST for the collapse functor and category structure,
  \item \textbf{Appendix P} for persistent homology collapse theory,
  \item \textbf{Appendix Q} for Ext-class collapse and gluing obstruction theory.
\end{itemize}

These elements allow us to treat the global regularity problem not merely as an analytic challenge, but as a structural resolution governed by collapse within a typed categorical space.

\subsection{1.6 Outline of This Work}

The rest of this manuscript is organized as follows:

\begin{itemize}
  \item Chapter 2 reviews Leray–Hopf weak solutions and classifies obstruction types.
  \item Chapter 3 introduces persistent homology for flow structures.
  \item Chapter 4 encodes the velocity field as a derived sheaf.
  \item Chapter 5 formulates the main collapse-based regularity theorem.
  \item Chapters 6–7 provide two independent proofs based on \( \mathrm{PH}_1 \) and \( \mathrm{Ext}^1 \) vanishing.
  \item Chapter 8 formalizes the collapse diagram and type-theoretic structure.
  \item Chapter 9 reinterprets legacy criteria from v5.3 in the collapse framework.
  \item Chapter 10 concludes with prospects for generalization to other PDEs.
\end{itemize}




\end{document}
